\section{Motivation und Problemstellung}
Klassische WordClocks stellen die Uhrzeit üblicherweise im Fünf-Minuten-Takt dar. 
Die Anzeige erfolgt in sprachlicher Form und verfolgt in vielen Fällen primär gestalterische Ziele. 
Der funktionale Umfang ist häufig auf die reine Zeitdarstellung beschränkt. 
Gleichzeitig steigt im technischen Umfeld der Anspruch an multifunktionale Systeme. 
Informationsdarstellungen werden zunehmend mit zusätzlichen Daten kombiniert. 
Hierzu zählen beispielsweise genaue Zeitinformationen, externe Sensordaten oder netzwerkbasierte Dienste. 

Zwischen dekorativ ausgelegten WordClocks und informationsorientierten Anzeigeeinheiten besteht eine konzeptionelle Differenz. 
Während klassische WordClocks eine reduzierte und sprachbasierte Zeitdarstellung bieten, verzichten viele multifunktionale Displays auf besondere oder ästhetische Darstellungsformen. 
Daraus ergibt sich eine technische Fragestellung: 

Es ist zu untersuchen, wie eine WordClock so konzipiert werden kann, dass die charakteristische wortbasierte Anzeige erhalten bleibt und gleichzeitig zusätzliche Funktionen technisch sinnvoll und den Designansprüchen entsprechend integriert werden. 
Dabei sind Anforderungen an Hardwarearchitektur, Softwarestruktur, Energieversorgung, Datenanbindung und Benutzerinteraktion zu berücksichtigen. 